\chapter{Pre-Trade Risk and Kill Switches}\label{ch:ll:pre-trade-risk-and-kill-switches}

On 1 August 2012, one of the eight servers of a broker-dealer's order router, Knight Capital Americas, still carried an old
piece of code that a new deployment should have replaced. For 212 incoming orders it sent child orders to the market ``without
regard to the number of share executions'' already received, in the words of the regulator's order: 4 million executions in 154
stocks, more than 397 million shares, in about 45 minutes, and a loss of \$460 million (One Quant Book~7, chapter~21, tells the
deployment story, and One Quant Book~11, chapter~27, the risk policy it produced). This chapter reads the same order as an
engineer would: as a list of requirements. Every paragraph that says what the firm lacked (a control that compared orders
leaving the router with those that entered it, a position limit that counted outstanding orders, a firm-wide threshold that
stopped orders automatically, a way to halt the router on its own misbehaviour) is a check to implement, on the path of every
order, in a few tens of nanoseconds. The chapter builds that gate, then replays a runaway of the 2012 shape through it with
one check at a time, and finds that most of the checks a firm would think of first would not have stopped it.

\section{The requirements document}

\begin{definition}[Risk gate, fail-closed design]\label{def:ll:pre-trade-risk-and-kill-switches:gate}
A \emph{risk gate}\index{risk gate} is the component on the path of every order, between the code that decides to send it and
the \omterm{def:ll:order-gateway-and-order-management:gateway}{order gateway}, that checks it against the firm's limits and refuses it if any is exceeded; nothing reaches the venue
without passing it. A \emph{fail-closed design}\index{fail-closed design} is one whose failures stop the activity it protects:
a risk gate that has lost its limits, its reference prices or its view of the firm's position refuses orders rather than
passing them.
\end{definition}

The regulator's order (Release 34-70694) is precise about what was missing. The firm had controls before orders reached the
router, but none after it: nothing compared ``orders leaving SMARS with those that entered it'', and there were no procedures
``to halt SMARS's operations in response to its own aberrant activity''. It had a price cap, 9.5\% from the best price when the
parent order arrived, which did not help while prices moved less than that and did not apply to orders meant for the opening
auction. It had position limits for some groups, but ``these limits did not account for the firm's exposure from outstanding
orders''. It had no firm-wide capital thresholds linked to automated controls, and its position monitor ``relied entirely on
human monitoring''. The rule it was charged under, the market access rule, asks for controls that prevent the entry of orders
that exceed ``pre-set credit or capital thresholds'' and of erroneous orders that ``exceed appropriate price or size
parameters, on an order-by-order basis or over a short period of time, or that indicate duplicative orders''. The engineering
translation is a list: a price collar, size and notional limits per order, a position limit that counts orders in flight, a
capital threshold across the firm, a rate limit, a duplicate check, and a kill switch that the gate itself can pull. The pre-trade
risk check and the price collar are defined in One Quant Book~11 (chapter~27), which sets the policy; this chapter is about
doing them fast and doing them right.

\begin{dated}{2026-09}{What the rules ask of pre-trade controls}\label{dat:ll:pre-trade-risk-and-kill-switches:rules}
In the United States, Rule 15c3-5 of the Securities Exchange Act (consulted September 2026) requires a broker or dealer with
market access to have controls that prevent the entry of orders exceeding pre-set credit or capital thresholds in the aggregate
for each customer and for itself, and of erroneous orders exceeding price or size parameters, order by order or over a short
period, or indicating duplicative orders. In the European Union, Commission Delegated Regulation (EU) 2017/589 requires an
investment firm to be able to cancel immediately any or all of its unexecuted orders on any or all venues (article 12); to have
price collars, maximum order values, maximum order volumes and maximum message limits (article 15), and repeated automated
execution throttles; and to generate real-time alerts within five seconds of the relevant event.
\end{dated}

The gate sits in the trading process, on the strategy engine's thread or right after it, before the \omterm{def:ll:order-gateway-and-order-management:gateway}{order gateway}
(\cref{fig:ll:pre-trade-risk-and-kill-switches:where}): the engine's order becomes the gateway's message only if the gate
says so. Everything else in the figure feeds it. The limits come from a risk service, as whole snapshots, with a \omterm{def:ll:protocols-i-fix:session}{heartbeat}; the
reference prices come from the \omterm{def:ll:the-order-book-builder:builder}{book builder}; the fills and cancellations come back from the gateway, so that the gate's view of
the position and of the orders in flight is the gateway's; and the kill switch can be pulled by the gate itself, by a person on
the risk desk, or by the venue.

\begin{omfigure}
\begin{tikzpicture}[font=\scriptsize,>=Stealth,box/.style={draw,rounded corners=2pt,minimum height=0.75cm,align=center,
  inner sep=3pt,fill=omDef!12}]
\node[box] (eng) at (0,0) {strategy\\ engine};
\node[box,fill=omThm!18,minimum width=2.2cm,minimum height=1.1cm] (gate) at (3.3,0) {\textbf{risk gate}\\ 12 checks, one mask};
\node[box] (gw) at (6.6,0) {order\\ gateway};
\node[box,fill=gray!10] (venue) at (9.4,0) {venue};
\node[box,fill=omMeth!15] (lim) at (1.6,2.0) {risk service:\\ limit snapshots,\\ heartbeats};
\node[box,fill=omMeth!15] (bb) at (5.0,2.0) {book builder:\\ reference prices};
\node[box,fill=omMeth!15] (desk) at (3.3,-2.1) {risk desk:\\ kill switch};
\draw[->,thick] ([yshift=2mm]eng.east) -- node[above]{order} ([yshift=2mm]gate.west);
\draw[->,thick] (gate) -- node[above]{accepted} (gw);
\draw[->,thick] (gw) -- (venue);
\node[box,fill=gray!10] (mon) at (9.4,2.0) {post-trade\\ monitor};
\draw[->] (lim) -- (gate);
\draw[->] (bb) -- (gate);
\draw[->] (desk) -- (gate);
\draw[->,dashed] (gw.south) -- ++(0,-0.7) -| node[pos=0.25,below]{fills, cancels} ([xshift=8mm]gate.south);
\draw[->,dashed] (venue) -- node[right]{drop copy} (mon);
\draw[->,omThm] ([yshift=-2.5mm]gate.west) -- node[below,text=omThm]{refused} ([yshift=-2.5mm]eng.east);
\end{tikzpicture}

\omcaption{Where the gate sits: on the path of every order, before the gateway, fed by limits with \omterm{def:ll:protocols-i-fix:session}{heartbeats}, reference
prices and the gateway's own view of fills and orders in flight. A post-trade monitor on the \omterm{def:ll:order-gateway-and-order-management:dropcopy}{drop copy} (chapter~21) is a
second line, not a substitute: it sees what has already happened.}\label{fig:ll:pre-trade-risk-and-kill-switches:where}
\end{omfigure}

\section{Checks in nanoseconds}

\begin{definition}[Worst-case exposure]\label{def:ll:pre-trade-risk-and-kill-switches:worst}
The \emph{worst-case exposure}\index{worst-case exposure} of a firm in an instrument is its position plus everything its open
orders could still add if they all filled: on the long side, the position plus the remaining quantity of every open buy order,
including those whose acknowledgement or cancellation has not arrived; on the short side, the same with sell orders. A limit is
respected in the worst case if a new order fits under it on top of this exposure.
\end{definition}

A \omterm{def:ll:pre-trade-risk-and-kill-switches:gate}{risk gate} is checked on every order, so it must cost little, and cost the same whether it accepts or refuses: a gate that is
fast on the common path and slow on the rare one is slowest exactly when a runaway is sending orders. The build's gate is
organised for that. The limits are held in dense tables indexed by instrument (the venue's locate code) and by strategy, not in
maps keyed by name. Whatever can be computed before the order arrives is: the collar is kept as two prices, recomputed when the
reference price or the limits change, so that the collar check is one comparison; the exposures are running sums, updated by
the reports as in chapter~21, so that the position check is an addition and a comparison. And every check is computed, without
a branch, into one bit of a mask (\cref{lst:ll:pre-trade-risk-and-kill-switches:check}); the order passes if the mask is zero,
and a refusal reports the lowest bit set, while the audit record keeps the whole mask, because the second reason matters when
the first is fixed.

\omcode{code/firm/riskgate/cpp/firm_riskgate.hpp}{124}{140}{Twelve checks, one mask: each condition sets one bit without a
branch, and a refusal reports the first.}\label{lst:ll:pre-trade-risk-and-kill-switches:check}

\Cref{fig:ll:pre-trade-risk-and-kill-switches:cost} measures it. With every order accepted and the duplicate check off, a check
costs about \qty{11}{\nano\second} at the median; the duplicate check, which compares the order with the strategy's sixteen
most recent ones, adds about \qty{10}{\nano\second}. The first version of that check compared five fields per remembered order
and cost more than all the other checks together; packing quantity and price into one word and instrument and side into another
(the wire's quantity and price are 32-bit fields, so the packing is exact) brought it to this. On the fixture's stream, where every kind of
refusal occurs and each check follows the parsing of its input line, a check costs about \qty{26}{\nano\second} at the median,
the same for accepted and refused orders: the mask makes the refusals no cheaper and no dearer.

\begin{omfigure}
\begin{tikzpicture}
\begin{axis}[width=11cm,height=5.4cm,ybar,bar width=14pt,ymin=0,ymax=45,xmin=-0.6,xmax=4.6,
  xtick={0,1,2,3,4},xticklabels={{fixture:\\ all},{fixture:\\ accepted},{fixture:\\ refused},{all accepted},{all accepted:\\ duplicates on}},
  xticklabel style={align=center,font=\footnotesize},ylabel={ns per check (median)},
  tick label style={font=\small},label style={font=\small},grid=major,grid style={gray!25},
  nodes near coords,every node near coord/.append style={font=\scriptsize},table/col sep=comma]
\addplot[fill=omThm!45,draw=omThm] table[x expr=\coordindex/2,y=ns,
  restrict expr to domain={mod(\coordindex,2)}{0:0}]{figdata/low-latency/22-pre-trade-risk-and-kill-switches/measured_check.csv};
\end{axis}
\end{tikzpicture}

\omcaption{The cost of one pass through the gate's twelve checks, median of TSC-timed calls: on the fixture's 4\,874 orders
(every kind of refusal, cold tables), split by outcome; and on a stream where every order passes, with the duplicate check off
and on. Measured on a laptop (Intel Core Ultra 7 155H) under WSL2. Data:
\texttt{bench\_risk.py}.}\label{fig:ll:pre-trade-risk-and-kill-switches:cost}
\end{omfigure}

Tens of nanoseconds are small next to the path of an order through the kernel and the network (chapter~1), and there is no
reason to make them smaller by checking less. The temptation runs the other way: to move the checks out of the trading
process, into a separate risk process that sees the orders on a ring and answers. That costs a round trip between cores, about \qty{200}{\nano\second} through the
rings of chapter~12, eight times the check itself, and, worse, it creates a second failure to handle: if the risk process stops answering, the trading process must
either wait or send unchecked. A gate in the trading process, with limits it can check alone, fails in one way only.

\section{Limits in a hierarchy, updated live}

Limits are set at several levels: the firm has a capital threshold, each desk a share of it, each strategy a rate, a number of
open orders and its own share, each instrument a collar, a size and a position limit (the risk hierarchy of One Quant Book~6,
chapter~29, and the policy of Book~11, chapter~27). An order is checked at every level it belongs to, and the gate keeps one
running exposure per level. The limits change during the day, when a trader asks for more or the risk desk takes some away,
and the gate applies changes the way the strategy engine applies parameters (chapter~20): as whole, versioned snapshots,
between two orders, never in the middle of a check, so that no order is checked against half an old snapshot and half a new
one.

The gate fails closed. A snapshot carries a maximum age: the risk service sends \omterm{def:ll:protocols-i-fix:session}{heartbeats}, and if the last snapshot or
\omterm{def:ll:protocols-i-fix:session}{heartbeat} is older than its age, every order is refused as stale. The build's fixture includes a risk service that goes
quiet for three seconds with a maximum age of two: from the moment its last \omterm{def:ll:protocols-i-fix:session}{heartbeat} is two seconds old until the next one
arrives, the gate refuses every order, 50 of them. The same holds for the reference prices behind the collar: an instrument whose reference is older than its maximum age
cannot be traded, because a collar around a stale price is no collar. Failing closed has a cost, since a strategy stops when a
monitoring service fails, and the firm must decide which services the trading depends on. The rule of the build is that
anything a check needs is such a service.

\section{Kill switches: block, cancel, flatten}

\begin{definition}[Mass cancel]\label{def:ll:pre-trade-risk-and-kill-switches:mass}
A \emph{mass cancel}\index{mass cancel} is a request that cancels many orders at once, all the open orders of a session, of an
instrument or of a side, either as one message the venue offers (the simulator's \texttt{M} message) or as a cancel for every
open order sent together; the orders remain in flight, and may still fill, until each cancellation is confirmed.
\end{definition}

A kill switch stops a node of the hierarchy (a strategy, a desk, the firm) and can do three things, in increasing order of
consequence. \emph{Block} refuses every new order under the node: nothing more goes out, and nothing already out is touched.
\emph{Cancel} also sends a \omterm{def:ll:pre-trade-risk-and-kill-switches:mass}{mass cancel} for the node's open orders; until each cancellation is confirmed, the orders count in
the exposure, and those that fill in the meantime are races of the kind of chapter~21. \emph{Flatten} also sends orders that
close the positions, at the reference price, which the collar admits; it is the most dangerous of the three, since it trades,
into a market the runaway may have moved, and the build returns those orders for a person to confirm rather than sending them.
Who pulls the switch is a policy question (Book~11): the gate pulls it on a breach of a hard limit, the risk desk can pull it
by hand, and many venues offer their own. Unkilling is always a person's decision.

\omcode{code/firm/riskgate/firm_riskgate.py}{198}{210}{The kill switch of the Python reference: block, cancel (a mass cancel of
the node's open orders), flatten (the closing orders, returned for confirmation).}\label{lst:ll:pre-trade-risk-and-kill-switches:kill}

\section{Where the checks sit}

Checks can sit in four places, and a firm uses several. In the trading process, as here: the cheapest and the only place
that sees every order before it leaves, with the firm's own view of the orders in flight. In a separate risk process or
appliance on the path: independent of the trading code, at the cost of a hop, and the only protection against a bug in
the trading process itself. At the broker, when the firm reaches the venue through a broker's sponsored access (One Quant
Book~1, chapter~4): the broker is responsible under the market access rule for the orders it lets through, whatever the
firm checks. And at the venue, whose collars, limits and kill switches protect the market rather than the firm. The 2012 router
had checks upstream and a monitor downstream, and nothing on the path of its own output. A useful test of any architecture is to
ask, for each component that generates orders, which check sees its output before the venue does.

\section{Tutorial: a runaway, one check at a time}\label{tut:ll:pre-trade-risk-and-kill-switches:tut}

\texttt{firm\_riskgate\_runaway.py} rebuilds the 2012 shape against Book 10's matching engine. The runaway sends buy child orders
of 100 shares in bursts of fifteen, a microsecond apart, every \qty{10}{\milli\second}: 1\,500 a second on average, the rate
of 4 million executions in 45 minutes, of about 99 shares each. It prices them five ticks above the last trade and never looks
at its fills. A liquidity provider keeps a ladder of sell orders and moves it up a tick every ten fills, so the runaway pushes
the price up as it buys. \Cref{tab:ll:pre-trade-risk-and-kill-switches:runaway} and
\cref{fig:ll:pre-trade-risk-and-kill-switches:runaway} give five seconds of it, with each check switched on in turn.

\begin{table}[ht]
\centering\small
\begin{tabular}{lrrrr}
\toprule
check switched on & orders sent & shares bought & first refusal & carried over 45 min \\
\midrule
none & 7\,500 & 750\,000 & --- & 405 million \\
collar, 5\% around the last trade & 7\,500 & 750\,000 & --- & 405 million \\
size and notional per order & 7\,500 & 750\,000 & --- & 405 million \\
throttle, 500 orders a second & 2\,545 & 254\,500 & \qty{40}{\milli\second} & 135 million \\
duplicates within \qty{100}{\milli\second} & 55 & 5\,500 & first repeat & 3.0 million \\
position 20\,000, fills only & 210 & 21\,000 & \qty{140}{\milli\second} & 21\,000 \\
position 20\,000, in flight & 200 & 20\,000 & \qty{130}{\milli\second} & 20\,000 \\
capital threshold, \$2.5 million & 249 & 24\,900 & \qty{160}{\milli\second} & 24\,900 \\
\bottomrule
\end{tabular}
\caption{A runaway of the 2012 shape, five seconds against the matching engine, with one check at a time. ``Carried over
45 minutes'' extends the rate of the last second when the runaway is still sending, and keeps the position when it has been
stopped. Data: \texttt{fig\_runaway.py} (deterministic).}\label{tab:ll:pre-trade-risk-and-kill-switches:runaway}
\end{table}

\begin{omfigure}
\begin{tikzpicture}
\begin{axis}[width=11cm,height=6cm,ymode=log,xmin=0,xmax=5000,ymin=50,ymax=2e6,
  xlabel={time since the runaway started (ms)},ylabel={position (shares)},
  tick label style={font=\small},label style={font=\small},grid=major,grid style={gray!25},
  legend pos=outer north east,legend style={font=\scriptsize,cells={anchor=west}},table/col sep=comma]
\addplot[omThm,thick] table[x=t_ms,y=none]{figdata/low-latency/22-pre-trade-risk-and-kill-switches/runaway_trace.csv};
\addlegendentry{none (or collar, size)}
\addplot[omMeth,thick] table[x=t_ms,y=throttle]{figdata/low-latency/22-pre-trade-risk-and-kill-switches/runaway_trace.csv};
\addlegendentry{throttle}
\addplot[gray,thick,dashed] table[x=t_ms,y=duplicates]{figdata/low-latency/22-pre-trade-risk-and-kill-switches/runaway_trace.csv};
\addlegendentry{duplicates}
\addplot[omDef,thick] table[x=t_ms,y=position_in_flight]{figdata/low-latency/22-pre-trade-risk-and-kill-switches/runaway_trace.csv};
\addlegendentry{position, in flight}
\addplot[black,thick,dotted] table[x=t_ms,y=capital_threshold]{figdata/low-latency/22-pre-trade-risk-and-kill-switches/runaway_trace.csv};
\addlegendentry{capital threshold}
\end{axis}
\end{tikzpicture}

\omcaption{The runaway's position, every \qty{50}{\milli\second}, with one check at a time. Checks on each order in isolation
never trip; a throttle slows the runaway and the duplicate check nearly stops it; a position limit that counts orders in flight,
or a capital threshold, stops it within 250 orders. Data: \texttt{fig\_runaway.py}.}\label{fig:ll:pre-trade-risk-and-kill-switches:runaway}
\end{omfigure}

The table is the chapter's argument. The checks that look at each order alone (collar, size, notional) never trip: every child
order is small, and priced near a market that the runaway itself is moving, so a collar around the last trade moves with it
(the price rose 7.5\% in the five seconds, more than the collar, and the collar never noticed). A throttle slows the runaway to
its rate and does not stop it: 135 million shares in 45 minutes is still a disaster. The duplicate check nearly stops this
runaway, because its child orders repeat, and would miss one that varied its quantity. What stops it is a limit on what the
firm holds, counted in the worst case: the position limit at 20\,000 shares after 200 orders, the capital threshold at
24\,900. Counting only fills overshoots by the orders in flight, here one burst of ten orders, 1\,000 shares; with more orders
in flight, or a slower report path, it overshoots by more. The kill switch changes what happens next: without it, the runaway
keeps sending and the gate keeps refusing, 7\,300 refusals in five seconds; with it, the strategy is blocked at the first
breach and its open orders are cancelled.

\section{Build: the risk gate}\label{bld:ll:pre-trade-risk-and-kill-switches:riskgate}

\begin{build}
\bfield{Purpose} The check on the path of every order, between the strategy engine (chapter~20) and the \omterm{def:ll:order-gateway-and-order-management:gateway}{order gateway}
(chapter~21); the policy it enforces is One Quant Book~11's (\texttt{firm.riskctl}, chapter~27), and its audit records go to the
binary log (chapter~23).
\bfield{Interface} Python reference \texttt{firm\_riskgate}: \texttt{Limits.from\_dict}, \texttt{from\_riskctl},
\texttt{Gate(limits, t)} with \texttt{check}, \texttt{set\_limits}, \texttt{heartbeat}, \texttt{set\_reference},
\texttt{on\_fill}, \texttt{on\_done}, \texttt{kill}, \texttt{unkill}, \texttt{audit}; \texttt{replay}, \texttt{decisions};
\texttt{firm\_riskgate\_runaway.run} against Book 10's engine. C++20 \texttt{firm::risk}: \texttt{parse\_limits},
\texttt{Gate} with the same operations and \texttt{Decision\{code, mask\}}. Rust \texttt{firm\_riskgate}: \texttt{Gate},
\texttt{parse\_limits}, \texttt{replay}.
\bfield{Rules} Every check computed into a mask, no branch on the outcome; limits as whole snapshots between orders; stale limits
or reference prices refuse (fail closed); exposure counts every order in flight; no allocation per check; a kill is undone
only by a person.
\bfield{Acceptance tests} \texttt{code/firm/riskgate/}: the fixture's 4\,874 orders replayed to the same decisions in Python, C++
and Rust, with every check refusing at least once; each check at its boundary; fail closed on stale limits and references; the
kill switch's three actions at firm, desk and strategy level; the running sums equal to a recount after random streams; zero
allocations; the runaway stopped by the position limit and the capital threshold, and not by the collar; Book 11's
\texttt{firm.riskctl} limits, exported for the gate, replayed in Python and C++.
\bfield{Stretch} A loss limit on realised and unrealised profit (exercise~7); per-venue message limits; limits shared by
several gateways through shared memory, with one writer; Book 11's monitors (loss, message rate) feeding the kill switch.
\end{build}

\begin{omsources}
\item U.S. Securities and Exchange Commission, Release No. 34-70694, \emph{In the Matter of Knight Capital Americas LLC}, 2013.
\item 17 CFR 240.15c3-5, \emph{Risk management controls for brokers or dealers with market access}.
\item Commission Delegated Regulation (EU) 2017/589 of 19 July 2016 (RTS 6).
\end{omsources}

\section{Exercises}

\begin{exercise}[$\star$]\label{exo:ll:pre-trade-risk-and-kill-switches:1}
The reference price is 25.0000 (250\,000 in units of \num{e-4}) and the collar is 500 basis points. Which of buy orders at
26.2500 and 26.2600, and sell orders at 23.7500 and 23.7400, pass?
\end{exercise}

\begin{exercise}[$\star$]\label{exo:ll:pre-trade-risk-and-kill-switches:2}
The long limit is 5\,000 shares. The position is 2\,000 long, and buy orders for 1\,500 and 1\,000 are open, one of them not
yet acknowledged. Is a buy of 600 accepted? What is the largest buy that is?
\end{exercise}

\begin{exercise}[$\star$]\label{exo:ll:pre-trade-risk-and-kill-switches:3}
From the regulator's totals, compute the incident's average number of executions a second and of shares per execution.
\end{exercise}

\begin{exercise}[$\star\star$]\label{exo:ll:pre-trade-risk-and-kill-switches:4}
A burst of fifteen identical child orders leaves a microsecond apart. The duplicate window is \qty{100}{\milli\second}. How
many of the fifteen pass, and how would the runaway have to change to pass them all?
\end{exercise}

\begin{exercise}[$\star\star$]\label{exo:ll:pre-trade-risk-and-kill-switches:5}
The limits' maximum age is two seconds and \omterm{def:ll:protocols-i-fix:session}{heartbeats} come every \qty{400}{\milli\second}. The risk service stops at time
$t$. Until when are orders accepted, and why is ``until the service comes back'' the wrong answer?
\end{exercise}

\begin{exercise}[$\star\star$]\label{exo:ll:pre-trade-risk-and-kill-switches:6}
A check costs about \qty{26}{\nano\second} on the fixture's stream. What fraction of a tick-to-trade budget of
\qty{1}{\micro\second} is that, and what would checking in a separate process on another core add?
\end{exercise}

\begin{exercise}[$\star\star\star$]\label{exo:ll:pre-trade-risk-and-kill-switches:7}
\emph{Coding.} Add a loss limit per strategy: realised profit from fills plus unrealised profit at the reference price, with
a kill (block and cancel) when it falls below the limit. Where is the loss updated, and why is it not a pre-trade check like the
others?
\end{exercise}

\begin{exercise}[$\star\star\star$]\label{exo:ll:pre-trade-risk-and-kill-switches:8}
\emph{Find the flaw.} ``Our risk checks run in their own process, which reads a copy of every order the gateway sends and sends
a kill as soon as it sees a breach. It keeps the checks out of the latency path.''
\end{exercise}

\section{Problem: Forty-Five Minutes}

\begin{problem}[{Weekend problem --- which checks would have stopped the 2012 runaway}]\label{pb:ll:pre-trade-risk-and-kill-switches:1}
Use \cref{tab:ll:pre-trade-risk-and-kill-switches:runaway} (\texttt{runaway.csv}) and the regulator's totals: 4 million
executions, more than 397 million shares, about 45 minutes.

\medskip\noindent\textbf{Part I --- The rates.}
\begin{enumerate}
\item Compute the incident's execution rate and check that the runaway's 1\,500 orders a second match it.
\item How many shares does the unchecked runaway buy in five seconds, and in 45 minutes at the same rate? Compare with the
incident.
\item In the simulation, the price rose by a tick every ten fills. By how much did it rise in five seconds?
\item Why did the collar never trip, although the price moved by more than the collar?
\end{enumerate}

\medskip\noindent\textbf{Part II --- The checks that do not stop it.}
\begin{enumerate}[resume]
\item Why do the size and notional limits per order never trip?
\item The throttle trips after \qty{40}{\milli\second}. Why so late, and what is the runaway's rate afterwards?
\item Carry the throttled rate over 45 minutes.
\item The duplicate check lets 55 orders through in five seconds. Which ones, and what small change to the runaway would defeat
it?
\end{enumerate}

\medskip\noindent\textbf{Part III --- The checks that do.}
\begin{enumerate}[resume]
\item How many orders does the position limit of 20\,000 shares, counted in flight, let through, and when does it trip?
\item The limit that counts only fills ends at 21\,000. Explain the difference.
\item In what circumstances would that difference be large?
\item The capital threshold is \$2.5 million. At about \$100 a share, why does it stop the runaway near 24\,900 shares?
\end{enumerate}

\medskip\noindent\textbf{Part IV --- The verdict.}
\begin{enumerate}[resume]
\item State the \emph{named result}: orders and shares before each check trips, carried over 45 minutes, against the
incident's totals.
\item What does the kill switch add to the position limit in this simulation?
\item The runaway traded 154 stocks. What does that change for a per-instrument position limit, and which check is immune to it?
\item Which of the missing controls listed in the regulator's order corresponds to each row of the table?
\item Why must the gate fail closed?
\item Why is a post-trade monitor, however fast, not a substitute for the gate?
\item What false positives would the duplicate check raise on an ordinary day?
\item In one sentence: what kind of limit stops a runaway?
\end{enumerate}
\end{problem}

\section{Interview questions}

\begin{interviewq}[$\star$ \iqroles{developer}]\label{iq:ll:pre-trade-risk-and-kill-switches:1}
List the pre-trade checks you would put on every order, and say which of them you would expect to catch a runaway algorithm.
\end{interviewq}

\begin{interviewq}[$\star\star$ \iqroles{developer}]\label{iq:ll:pre-trade-risk-and-kill-switches:2}
How do you make a risk check cost a few nanoseconds, and cost the same whether it accepts or refuses?
\end{interviewq}

\begin{interviewq}[$\star\star$ \iqroles{developer, risk}]\label{iq:ll:pre-trade-risk-and-kill-switches:3}
Your risk service stops sending limits. What does the trading system do, and why?
\end{interviewq}

\begin{interviewq}[$\star\star$ \iqroles{developer, trader}]\label{iq:ll:pre-trade-risk-and-kill-switches:4}
What should a kill switch do, and who should be able to pull it and undo it?
\end{interviewq}

\begin{interviewq}[$\star\star$ \iqroles{developer}]\label{iq:ll:pre-trade-risk-and-kill-switches:5}
How do you change a strategy's limits during the day without stopping it?
\end{interviewq}

\begin{interviewq}[$\star\star\star$ \iqroles{developer, risk}]\label{iq:ll:pre-trade-risk-and-kill-switches:6}
Read the 2012 incident as a requirements document. Design the controls, where they sit, and how you would test that they stop a
runaway.
\end{interviewq}
